\begin{problem}{맛있는 스콘 만들기}{standard input}{standard output}{3 seconds}{1024 megabytes}

도현이는 SCON 참가자들에게 나눠줄 간식으로 스콘을 만들기로 했다.  

도현이가 만들려는 스콘은 1초마다 오븐의 온도를 조절해 가며 만드는 특별한 스콘이며, 총 $N$초 동안 구워야 한다.  

모든 시각 $t \ (1 \leq t)$마다 최적의 온도 $k_t$가 존재하고, $t-1$초부터 $t$초까지 오븐의 온도를 $k$로 유지했다면 스콘의 맛은 $M - |k_t - k|$ 만큼 올라간다.  

스콘을 만드는데 사용할 오븐은 $1$부터 $M$까지의 정수 온도로만 조절이 가능하다.  
처음에는 자유롭게 오븐의 온도를 설정할 수 있지만, 이후에는 $C$의 정수 배만큼만 온도를 높이거나 낮출 수 있다. 또, 한 번에 $D$를 초과해서 온도를 높이거나 낮출 수 없다. 예를 들어, 현재 오븐의 온도가 182이고, $C=5, D=10$일 경우, 조절 가능한 온도는 $172, 177, 182, 187, 192$뿐이다.  

처음 스콘의 맛을 $0$이라고 했을 때, 매초 온도를 적절히 조절해서 스콘을 만들었을 때 스콘의 맛의 최댓값을 구해보자. 처음 스콘을 오븐에 넣었을 때의 시간이 $0$초이며, 이때 오븐의 온도는 자유롭게 설정할 수 있다.  

\InputFile
첫 번째 줄에 네 정수 $N, M, C, D$가 공백으로 구분되어 주어진다. $(1 \leq N \leq 1000, 1 \leq M \leq 5000, 1 \leq C \leq D \leq M, D \equiv 0 \pmod{C})$

두 번째 줄에 정수 $k_1, k_2, ... , k_N$이 주어진다. $(1 \leq k_t \leq M)$


\OutputFile
첫 번째 줄에 $N$초 후 스콘의 맛의 최댓값을 출력한다.

\Examples

\begin{example}
\exmpfile{example.01}{example.01.a}%
\exmpfile{example.02}{example.02.a}%
\end{example}

\end{problem}

